\section{Problem Formulation}
\label{sec:problem}

\subsection{Notation}

Consider an inner equi-join $F \bowtie_{\kappa} D$ on key $\kappa$, where $F$ is
the larger, skewed side with $N$ rows and $D$ the probe side with $M$ rows.
Spark hash-partitions both sides into $p$ shuffle partitions
(\texttt{spark.sql.shuffle.partitions}) and executes a sort-merge join with $C$
concurrently schedulable task slots.

Let the key frequencies on $F$ follow a truncated Zipf law over $n$ distinct
keys, $\Pr[\kappa = i] \propto i^{-\theta}$, and let key $1$ be the hot key.
Write

\begin{itemize}
\item $H$ for the rows of $F$ carrying the hot key,
\item $P$ for the rows of $D$ carrying the hot key.
\end{itemize}

Both are obtainable before execution, from table statistics, a sampled
frequency sketch, or the catalogue.

It is convenient to express $P$ relative to the probe side's mean rows per key,
$\bar{m} = M/n$, through a \emph{hot-key multiplier}
\begin{equation}
\mu \;=\; \frac{P}{\bar{m}} ,
\label{eq:mu}
\end{equation}
so that $\mu = 1$ is a probe side with no concentration at all --- one-sided
skew --- and larger $\mu$ makes the skew progressively two-sided. Treating the
degree of two-sidedness as continuous rather than binary matters for two
reasons, one practical and one methodological, and we return to both in
Section~\ref{sec:methodology}.

\subsection{The skew ratio}

The quantity that governs everything is dimensionless:

\begin{equation}
\rho \;=\; \frac{H}{N/p}
\label{eq:rho}
\end{equation}

the hot key's rows divided by the mean rows per shuffle partition. It says how
many ordinary partitions' worth of work has been concentrated into one.

Equation~\eqref{eq:rho} is a property of the data and the partition count
alone. It contains no hardware term, which is what allows a rule built on it to
be stated once and applied on any cluster. When $\rho \le 1$ there is no
straggler to remove. When $\rho \gg 1$ the stage's completion time is
essentially the hot task's completion time, regardless of how many executors are
provisioned.

\subsection{Why $\mu$ cannot be taken to its limit}

There is a natural temptation to construct the two-sided case by drawing both
sides from the same Zipf law. It does not work, and the reason is worth stating
because it constrains any experiment in this area.

If both sides follow the same law, the hot key contributes $H \times P$ rows to
the join output on its own. At the scales used here that is order $10^{9}$
rows from a single key, a thousandfold amplification of the fact table --- the join has become a near-cartesian product, and
any measurement of it reports the cost of materialising that product rather
than the cost of the skew mitigation under test. Bounding $\mu$ explicitly
keeps the output tractable and, more usefully, promotes $P$ to an independent
variable. Equation~\eqref{eq:kstar} predicts that the optimal salt cardinality
falls as $1/\sqrt{P}$; that prediction is only testable if $P$ can be swept.

\subsection{Two mechanisms, one problem}

\subsubsection{AQE skew-join}
Spark's \texttt{OptimizeSkewedJoin} rule classifies a shuffle partition as
skewed when its size exceeds \emph{both} an absolute threshold
(\texttt{skewedPartitionThresholdInBytes}) and a multiple
(\texttt{skewedPartitionFactor}) of the median partition size. It then splits
the offending partition into sub-partitions of roughly the advisory size and
\emph{replicates the corresponding partition of the other side} to each split,
so that every sub-partition retains its join partners~\cite{spark_aqe_docs,
xue2024adaptive}.

The replication is the hinge. If the other side's matching partition is small
--- the one-sided case --- replication is nearly free and the rule works as
advertised. If the other side is also concentrated on the same key --- the
two-sided case --- then each of the $s$ splits inherits a large partner, the
probe-side work multiplies by $s$, and the relief on the critical path is
partly or wholly cancelled. Metwally~\cite{metwally2025amjoin} identifies
doubly-skewed keys as the hard case for exactly this reason.

\subsubsection{Salting}
Salting rewrites the key. On $F$, the hot key is replaced by
$(\kappa, \sigma)$ with $\sigma$ drawn uniformly from $\{0,\dots,k-1\}$; on
$D$, each matching row is emitted $k$ times, once per salt value. One heavy
key becomes $k$ lighter ones.

We distinguish two variants, because their costs differ by orders of magnitude
and the literature does not separate them:

\begin{description}
\item[Uniform salting] every row of $D$ is replicated $k$ times. This is the
form shown in most tutorials. Replication cost scales with $M$.
\item[Selective salting] only rows whose key is hot are replicated; all other
rows carry salt $0$ on both sides and join exactly as before. Replication cost
scales with $P$.
\end{description}

Since $P \ll M$ in any realistic fact--dimension join, the distinction is not a
detail. We evaluate both.

\subsection{What is actually being traded}

Salting does not remove work; it moves it. Splitting the hot key across $k$
sub-keys reduces the critical task's input from $H$ to approximately $H/k$
rows, and simultaneously introduces $P(k-1)$ probe-side rows that did not
previously exist. One term falls, the other rises, and the practitioner is
choosing a point on that curve without being told there is a curve.

The next section prices it.
